\documentclass[11pt]{article}

\usepackage[utf8]{inputenc}
\usepackage[margin=2.5cm]{geometry}
\usepackage{amsmath}
\usepackage{amssymb}
\usepackage{enumitem}
\usepackage{xcolor}
\usepackage[colorlinks=true,linkcolor=blue,urlcolor=blue]{hyperref}

% --- helpers --------------------------------------------------------------
% Inline file / identifier reference (handles underscores safely).
\newcommand{\file}[1]{\texttt{\detokenize{#1}}}

% A puzzle entry: fields are filled with the \field macro below.
\newcommand{\field}[2]{\smallskip\noindent\textbf{#1:}\ #2\par}

% Status badges.
\newcommand{\statusok}{\textcolor{green!55!black}{\textbf{working}}}
\newcommand{\statusstub}{\textcolor{orange!85!black}{\textbf{stub}}}
\newcommand{\statuswip}{\textcolor{red!75!black}{\textbf{work in progress}}}

\title{\textbf{Math Puzzles}\\[2pt]\large Problem \& Solution Index}
\author{Ruben Colomina}
\date{\today}

\begin{document}
\maketitle

\begin{abstract}
\noindent
This document indexes the recreational mathematics, algorithm, and
cryptography puzzles collected in the \file{math_puzzles} repository. Each
topic maps to a self-contained directory; every puzzle below records its
problem statement, the idea behind the solution, its computational cost, the
source files involved, and its current status. See
Section~\ref{sec:intro} for how the collection is organised, built, and
extended.
\end{abstract}

\tableofcontents
\newpage

% =========================================================================
\section{Introduction and Conventions}\label{sec:intro}

The repository is a personal collection of puzzles grouped \emph{by topic}
rather than by language. Solutions are written in C++, Python, and Rust
depending on the problem. Many originate from Codility lessons and from the
book \emph{Understanding Cryptography}.

\subsection*{Layout}
Each top-level directory is an independent topic. Shared tooling lives in
\file{common/}:
\begin{itemize}[nosep]
  \item \file{common/log.h} --- header-only leveled logging for C++.
  \item \file{common/log.py} --- matching logging helper for Python.
  \item \file{common/topic.mk} --- shared Make rules every topic Makefile includes.
\end{itemize}

\subsection*{Building and running}
A root \file{Makefile} auto-discovers every topic directory that includes
\file{common/topic.mk}:
\begin{itemize}[nosep]
  \item \texttt{make} or \texttt{make build} --- compile every C++ topic.
  \item \texttt{make run} --- build, then run every solution.
  \item \texttt{make clean} --- remove built binaries.
  \item \texttt{make docs} --- rebuild this PDF.
\end{itemize}
Every \file{*.cpp} in a topic builds to a binary of the same base name, so a
new source file is picked up with no Makefile edits.

\subsection*{Logging}
Diagnostic output is routed through the shared loggers and gated by the
\texttt{VERBOSE} environment variable (levels: \texttt{0}=error,
\texttt{1}=warn, \texttt{2}=info [default], \texttt{3}=debug). Final answers
are logged at \emph{info}; step-by-step traces at \emph{debug}:
\begin{center}
\texttt{./odd\_occur\_array} \quad vs. \quad \texttt{VERBOSE=3 ./odd\_occur\_array}
\end{center}

\subsection*{Adding a new puzzle}
The helper \file{scripts/new_puzzle.sh} scaffolds a new puzzle already wired
into the build and logging conventions:
\begin{center}
\texttt{scripts/new\_puzzle.sh <topic> <name> [cpp|py]}
\end{center}
For a C++ puzzle it stamps out a logging-ready source file and, for a new
topic, a one-line Makefile that inherits the shared rules --- after which the
root build discovers it automatically.

% =========================================================================
\section{Arrays Handling}
Directory \file{arrays_handling/} --- classic linear-scan array puzzles.

\subsection{Minimum split difference}
\field{Problem}{Given an array $A$ of $N$ integers, split it at some index
$p$ ($1 \le p < N$) into a left part $A[0..p-1]$ and a right part
$A[p..N-1]$. Minimise the absolute difference
$\bigl|\sum_{\text{left}} - \sum_{\text{right}}\bigr|$.}
\field{Approach}{Compute the total right sum once, then sweep $p$ from left to
right, moving one element from the right sum into the left sum at each step and
tracking the minimum absolute difference. This turns a naive
$O(N^2)$ recomputation into a single pass.}
\field{Complexity}{$O(N)$ time, $O(1)$ extra space.}
\field{Files}{\file{min_split_diff.cpp}}
\field{Status}{\statusok}

\subsection{Missing element}
\field{Problem}{An array contains $N$ distinct integers taken from
$\{1,\dots,N+1\}$ with exactly one value missing. Find the missing value.}
\field{Approach}{The sum of $1,\dots,N+1$ is $\tfrac{(N+1)(N+2)}{2}$. Subtract
the actual array sum from this closed-form total; the difference is the
missing element. Wide integer types guard against overflow for large $N$.}
\field{Complexity}{$O(N)$ time, $O(1)$ extra space.}
\field{Files}{\file{missing_element.cpp}}
\field{Status}{\statusok\ (the sum method is correct; the bundled test driver
feeds values $2,\dots,N+1$ rather than $1,\dots,N+1$, so its printed number is
illustrative of the mechanics rather than a clean demo)}

\subsection{Odd occurrence in an array}
\field{Problem}{Every value in the array appears an even number of times
except one, which appears an odd number of times (in the unpaired variant,
exactly once). Find that value.}
\field{Approach}{Two solutions are provided: a map-based count that returns the
value seen an odd number of times, and the classic XOR fold --- XOR-ing all
elements cancels paired values ($x \oplus x = 0$), leaving the lone value.}
\field{Complexity}{Map: $O(N\log N)$ / $O(N)$ space. XOR: $O(N)$ time,
$O(1)$ space.}
\field{Files}{\file{odd_occur_array.cpp}}
\field{Status}{\statusok}

% =========================================================================
\section{Counting Elements}
Directory \file{counting_elements/}.

\subsection{Smallest missing positive integer}
\field{Problem}{Given an array of integers (possibly negative or duplicated),
find the smallest positive integer that does \emph{not} occur in it. For
example $[1,3,6,4,1,2] \to 5$ and $[-1,-3] \to 1$.}
\field{Approach}{Insert the values into an ordered set, discard everything
$\le 0$, and walk the remaining positive values looking for the first gap:
if the adjacent difference exceeds $1$ (or the sequence does not start at
$1$), the missing integer sits in that gap.}
\field{Complexity}{$O(N\log N)$ time (set construction).}
\field{Files}{\file{missingint.cpp}}
\field{Status}{\statusok}

% =========================================================================
\section{Stacks and Queues}
Directory \file{stack_queues/}.

\subsection{Fish competition}
\field{Problem}{$N$ fish move along a river; each has a size $A[i]$ and a
direction $B[i]$ ($1$ = downstream, $0$ = upstream). When a downstream fish
meets an upstream one they fight and the larger eats the smaller. How many
fish survive?}
\field{Approach}{Scan the fish left to right maintaining a stack of survivors.
A downstream fish on top of the stack meeting an upstream incoming fish
triggers a fight loop that pops eaten fish; otherwise the fish is pushed.
Fish that never meet an opponent survive automatically.}
\field{Complexity}{$O(N)$ time, $O(N)$ space.}
\field{Files}{\file{fish_compet.cpp}}
\field{Status}{\statusok}

\subsection{Properly nested brackets}
\field{Problem}{Given a string of the brackets \texttt{()[]\{\}}, decide
whether it is properly nested.}
\field{Approach}{The iterative solution pushes opening brackets onto a stack
and pops on a matching closer; the string is well-formed iff the stack ends
empty. A recursive variant is also sketched in the file.}
\field{Complexity}{Iterative: $O(N)$ time, $O(N)$ space.}
\field{Files}{\file{properly_nested.cpp}}
\field{Status}{\statusok\ (iterative solver; the recursive variant is
experimental)}

% =========================================================================
\section{Graph and Tree Algorithms}
Directory \file{graph_algorithms/}.

\subsection{Visible nodes in a binary tree}
\field{Problem}{A node in a binary tree is \emph{visible} if no node on the
path from the root to it holds a strictly greater value. Count the visible
nodes.}
\field{Approach}{An explicit stack performs a depth-first traversal carrying
the maximum value seen so far along each root-to-node path; a node is counted
when its value is at least that running maximum. The \file{binary_tree.h}
header defines the \texttt{TreeNode} type together with DFS printing and the
visibility count.}
\field{Complexity}{$O(N)$ time, $O(h)$ stack space for height $h$.}
\field{Files}{\file{visible_nodes_bin_tree.cpp}, \file{binary_tree.h}}
\field{Status}{\statusok}

\subsection{Word machine (reverse-polish stack calculator)}
\field{Problem}{Evaluate a whitespace-separated program of integers and the
operators \texttt{POP}, \texttt{DUP}, \texttt{+}, \texttt{-} against a stack
machine, returning the value left on top.}
\field{Approach}{Tokenise the input and fold it over a stack: numbers are
pushed, operators consume and push results. The intended algorithm is present
but the driver still short-circuits and does not compile cleanly, so it is
excluded from the default build.}
\field{Complexity}{$O(T)$ time in the number of tokens $T$ (target).}
\field{Files}{\file{word_machine.cpp}}
\field{Status}{\statuswip}

% =========================================================================
\section{Codility Tests}
Directory \file{codility_test/}.

\subsection{Allocate families}
\field{Problem}{An aircraft has $N$ rows of seats laid out
\texttt{ABC~DEFG~HJK}. Given the already-reserved seats, count how many
3-person families can be seated together (each family in three adjacent seats
within one of the blocks \texttt{ABC}, \texttt{DEF}/\texttt{EFG}, or
\texttt{HJK}).}
\field{Approach}{Parse the reserved seats into a per-row occupancy grid, then
for each row check which of the three seat blocks remain fully free and tally
the achievable family placements.}
\field{Complexity}{$O(N + R)$ time for $R$ reservations.}
\field{Files}{\file{allocate_families.cpp}}
\field{Status}{\statusok}

\subsection{Weeks in Hawaii}
\field{Problem}{Placeholder for a date/interval Codility task.}
\field{Approach}{Not yet implemented --- the solution body is empty. Retained
as a stub so the intended task is tracked.}
\field{Complexity}{---}
\field{Files}{\file{weeks_in_hawaii.cpp}}
\field{Status}{\statusstub}

% =========================================================================
\section{Combinatorics}
Directory \file{combinatory/}.

\subsection{Cutting strings}
\field{Problem}{Enumerate every way to insert dash separators between the
characters of a string, e.g. \texttt{"12345"} yields
\texttt{"12345"}, \texttt{"1234-5"}, \dots, \texttt{"1-2-3-4-5"}.}
\field{Approach}{Each of the $n-1$ gaps between characters is independently
either a cut or not, giving $2^{\,n-1}$ splits. \texttt{itertools.product}
enumerates the binary choice vectors and each is rendered into a dashed
string.}
\field{Complexity}{$O(n \cdot 2^{\,n-1})$ to materialise all splits.}
\field{Files}{\file{cutting_strings.py}}
\field{Status}{\statusok}

% =========================================================================
\section{Parallel Computing}
Directory \file{parallel_compu_primes/}.

\subsection{Parallel prime sieve (OpenMP)}
\field{Problem}{Determine primality for every integer below
$10^{6}$, using multiple CPU threads.}
\field{Approach}{An OpenMP \texttt{parallel for} distributes trial-division
primality tests across threads, filling a boolean array; the primes just below
the ceiling are then reported. Compiled with \texttt{-fopenmp} via the topic
Makefile.}
\field{Complexity}{$O\!\left(\sum_{n} \sqrt{n}\right)$ work, divided across
threads.}
\field{Files}{\file{primes.cpp}}
\field{Status}{\statusok}

% =========================================================================
\section{Cryptography}
Directory \file{crypto/}.

\subsection{Special-sum cryptarithm}
\field{Problem}{Find distinct positive digits $A,B,C$ such that the addition
\[
\underline{AAA} + \underline{BBB} + \underline{CCC} = \underline{BAAC}
\]
holds, where each letter denotes a repeated digit.}
\field{Approach}{The equation reduces to
$111(A+B+C) = 1000B + 110A + C$. A brute-force search over all distinct digit
triples $(A,B,C)$ checks this identity; the solution is
$(A,B,C) = (9,1,8)$.}
\field{Complexity}{$O(9^3)$ constant-bounded search.}
\field{Files}{\file{special_sum/special_sum.py}}
\field{Status}{\statusok}

\subsection{Substitution cipher --- frequency analysis}
\field{Problem}{\emph{Understanding Cryptography}, Exercise 1.1: recover the
plaintext of a mono-alphabetic substitution cipher without the key.}
\field{Approach}{Compute the relative frequency of each ciphertext letter and
align it, in sorted order, against the known frequency distribution of English
letters to seed a decryption map. A few manual transpositions, guided by the
partially readable output, complete the mapping and reveal the plaintext
(a passage on Okinawan karate).}
\field{Complexity}{$O(L)$ for text length $L$, plus a small manual refinement.}
\field{Files}{\file{undertanding_cryptography/chapter01/problem01.py}}
\field{Status}{\statusok}

\subsection{ElGamal encryption (Rust)}
\field{Problem}{Demonstrate ElGamal public-key encryption over an elliptic
curve: generate a key pair, encrypt a message point, and verify that
decryption recovers it.}
\field{Approach}{A Cargo project using the \texttt{rust-elgamal} and
\texttt{rand} crates builds a decryption key, derives the encryption key,
encrypts a scalar multiple of the generator, and asserts that decryption
returns the original point.}
\field{Complexity}{Group-operation bound of the underlying curve arithmetic.}
\field{Files}{\file{elgamal-rust/src/main.rs}}
\field{Status}{\statusok\ (built with \texttt{cargo build} / \texttt{cargo run})}

% =========================================================================
\section{Miscellaneous}

\subsection{The integral of $(-1)^x$ over $[0,1]$}
\field{Problem}{Evaluate $\displaystyle\int_0^1 (-1)^x\,\mathrm{d}x$, where the
integrand takes complex values.}
\field{Approach}{Existence follows because $|(-1)^x| \le 1$ is integrable.
Writing $(-1)^x = e^{i\pi x}$ via Euler's identity gives
\[
\int_0^1 e^{i\pi x}\,\mathrm{d}x
  = \left.\frac{e^{i\pi x}}{i\pi}\right|_0^1
  = \frac{2i}{\pi},
\]
whose modulus $\tfrac{2}{\pi} < 1$ matches the bound. A full write-up is in
\file{integrate/amazing_integral.tex}.}
\field{Complexity}{---}
\field{Files}{\file{integrate/amazing_integral.tex}}
\field{Status}{\statusok\ (LaTeX write-up)}

\subsection{Alea --- a threaded pseudo-random generator}
\field{Problem}{Provide a linear-congruential pseudo-random number generator
that can be driven either serially or by a background producer thread feeding
a shared queue.}
\field{Approach}{The header \file{random_numbers/alea.h} declares an
\texttt{Alea} class wrapping LCG constants and a circular buffer / deque, with
a \texttt{pthread}-based producer for the parallel path.}
\field{Complexity}{$O(1)$ amortised per generated number.}
\field{Files}{\file{random_numbers/alea.h}}
\field{Status}{\statusok\ (header-only helper)}

\end{document}
